\section{Discussion}\label{sec:discussion}

\subsection{Interpretation of results}\label{subsec:interpretation}

Our findings reveal a nuanced picture of how proof decomposition affects the
optimization landscape. We discuss three aspects in turn.

\para{Local flatness via spectral orthogonality.}
Theorem~\ref{thm:hessian-bound} identifies spectral orthogonality of per-step
Hessians as the mechanism by which CoT achieves flatter minima. This is not
merely a consequence of having more parameters or longer sequences; it follows
from the structural independence of different reasoning steps. When modus ponens
and conjunction elimination involve different parameter subspaces, their
Hessians contribute curvature in orthogonal directions, and the top eigenvalue
is controlled by the worst single step. This spectral decoupling is a geometric
consequence of the modular structure of logical reasoning.

\para{The flat-but-isolated phenomenon.}
The most surprising finding is that CoT solutions are individually flatter but
globally more isolated. Proposition~\ref{prop:barrier-super} explains the
mechanism: loss barriers accumulate across steps. Each intermediate prediction
must interpolate correctly for the total barrier to remain low, and this
becomes harder as $K$ increases. This reveals an important distinction between
two senses of ``landscape smoothness'': local curvature (captured by Hessian
eigenvalues) and global connectivity (captured by loss barriers).

An alternative explanation for the higher CoT barriers involves the
specialization of intermediate representations. If each trained CoT model
develops slightly different internal derivation strategies---for instance,
applying rules in different orders---then interpolating between such models
disrupts the internal coherence of the reasoning chain, even if each
individual model sits in a flat basin.

\para{Cross-type barriers.}
The lowest barriers occur between direct and CoT models trained with different
seeds ($0.73 \pm 0.01$), rather than between same-type models. This suggests
that the proof generation strategy determines which \emph{region} of parameter
space is explored, and these regions are closer to each other than to other
solutions of the same type. One interpretation is that direct and CoT models
share a common ``logical core'' in parameter space, with the CoT model adding
specialized parameters for intermediate step generation that do not strongly
conflict with the direct model's parameters.

\subsection{Connections to related work}\label{subsec:related}

\para{Loss landscape geometry.}
Li et al.~\cite{li2018visualizing} showed that architectural choices (skip
connections, width) fundamentally alter landscape geometry. Our work extends
this finding to the \emph{output strategy} dimension: even with identical
architectures, the choice of how to structure the target sequence reshapes the
landscape. The filter-normalized visualization technique of~\cite{li2018visualizing}
was essential for obtaining meaningful cross-strategy comparisons.

Dinh et al.~\cite{dinh2017sharp} warned that sharpness measures can be
manipulated by reparametrization in networks with positively homogeneous
activations. Our models use GELU activations, which are not positively
homogeneous, partially mitigating this concern. Nevertheless, we used
filter normalization for 2D surface visualization and report multiple
complementary sharpness measures to ensure robustness.

\para{Sharpness and generalization.}
Foret et al.~\cite{foret2021sharpness} connected sharpness to generalization
via PAC-Bayes bounds. Our CoT models have both lower sharpness and lower
validation loss, consistent with this theory. An interesting direction is
whether training direct models with SAM would close the generalization gap by
explicitly seeking flatter regions, testing whether the landscape geometry is
a \emph{cause} or merely a \emph{correlate} of the performance difference.

\para{Mode connectivity.}
Ainsworth et al.~\cite{ainsworth2023git} showed that independently trained
image classifiers lie in a single basin after permutation alignment.  Our
finding that CoT models have \emph{higher} same-type barriers (without
permutation alignment) suggests either that the single-basin phenomenon does
not extend to autoregressive sequence models, or that permutation alignment is
essential for revealing the underlying connectivity. Applying Git
Re-Basin~\cite{ainsworth2023git} to transformer-based theorem provers is an
important direction for future work.

\para{Neural theorem proving.}
The step-by-step approach of GPT-f~\cite{polu2020generative} and
HTPS~\cite{lample2022hypertree} provides dense per-step supervision analogous
to our CoT format, while Baldur~\cite{first2023baldur} represents the direct
approach. Our landscape analysis provides a geometric explanation for the
empirical observation that step-by-step methods tend to converge more reliably.
DeepSeek-Prover-V2~\cite{ren2025deepseek} introduces consistency rewards that
explicitly couple CoT structure to formal proofs; this can be interpreted as a
landscape-shaping mechanism that further enhances the spectral orthogonality
identified in Theorem~\ref{thm:hessian-bound}.

\subsection{Limitations}\label{subsec:limitations}

Several limitations constrain the scope of our conclusions.

\para{Scale.} Our models have 106{,}027 parameters, while state-of-the-art
theorem provers have billions. Loss landscape geometry may change qualitatively
at larger scale. In particular, the spectral orthogonality assumption in
Theorem~\ref{thm:hessian-bound}(ii) may break down in overparameterized
regimes where per-step Hessians share significant spectral overlap.

\para{Task complexity.} Propositional logic is decidable and admits short
proofs. Real theorem proving in dependent type theories (Lean~4, Isabelle)
involves complex term construction and genuine mathematical reasoning. Our
results establish the existence of geometric differences in a controlled
setting but do not guarantee they persist in more complex domains.

\para{Statistical power.} Three runs per condition provide limited statistical
power. The consistent directional effects are encouraging, but the magnitude of
the differences and their statistical significance would benefit from
replication with 10 or more seeds.

\para{Confounding factors.} CoT outputs are longer than direct outputs.
The models therefore process different amounts of target sequence, confounding
the generation strategy effect with a sequence length effect. A more careful
controlled experiment might fix the total output length and compare the effect
of decomposition structure alone.

\para{Reparametrization.} While filter normalization addresses some concerns
raised by Dinh et al.~\cite{dinh2017sharp}, the Hessian eigenvalue and SAM
sharpness measures are not fully reparametrization-invariant. Developing
invariant measures suitable for autoregressive models remains an open problem.
